\section{Concrete Improvements Supplied by the \(S^3\) Geometry}
The role of the \(S^3\) construction should be stated conservatively.  Its main
achievement so far is not a proof of the missing positivity statement, but a
sequence of reductions showing that the obstruction seen in the
Connes--Consani--Moscovici/CCM framework is much smaller and more structured
than the original formulation suggests.
\subsection{Exact removal of one exceptional scalar mode}
\begin{proposition}[Exact zero-mode cancellation]
For the distinguished test vector \(h_\lambda\) arising from the geometric
construction,
\[
    \widehat h_\lambda(0)=0.
\]
Thus one of the two scalar endpoint obstructions occurring in the CCM
finite-dimensional approximation disappears identically.
\end{proposition}
\paragraph{Proof sketch.}
The \(S^3\) model places \(h_\lambda\) in a geometrically distinguished
combination of spherical modes.  Fourier evaluation at the trivial character
is the corresponding scalar projection.  The coefficients forced by the
\(S^3\) symmetry have opposite contributions in this projection, and hence
cancel exactly.  This is an algebraic consequence of the mode decomposition,
rather than an asymptotic estimate.
\paragraph{Relation to the Connes gap.}
Connes' positivity strategy requires controlling the discrepancy introduced
when the global Weil form is compressed to the finite/prolate approximation.
A priori there are exceptional scalar directions not controlled by the
bulk spectral approximation.  The proposition eliminates one such direction
completely.  Hence this part of the positivity defect is not a genuine
obstruction.
\subsection{The surviving endpoint defect is forced to the tunnelling scale}
\begin{proposition}[Suppression of the second scalar defect]
The remaining endpoint value satisfies
\[
    h_\lambda(0)=O(1-\chi_4),
\]
where \(\chi_4\) is the relevant prolate concentration eigenvalue.
\end{proposition}
\paragraph{Proof sketch.}
Writing the geometrically selected vector in the \(S^3\)-adapted low-mode
basis, the leading contribution which could produce an \(O(1)\) endpoint
term cancels.  What remains is multiplied by the failure of the fourth mode
to be perfectly concentrated.  This failure is exactly \(1-\chi_4\).
\paragraph{Relation to the Connes gap.}
The exceptional contribution therefore lives on the same exponentially
small scale as the prolate leakage itself.  Instead of separately controlling
an arbitrary endpoint error, the Connes positivity problem is reduced to
understanding a single tunnelling scale
\[
    \varepsilon_\lambda:=1-\chi_4.
\]
This is important because it makes the exceptional and bulk approximation
errors commensurable.
\subsection{Exact geometric formula for Fourier leakage}
\begin{proposition}[Mode-by-mode leakage identity]
If the \(S^3\)-selected vector has the relevant decomposition
\[
    h_\lambda=a_\lambda h_0+b_\lambda h_4,
\]
then
\[
 \bigl\|(1-\widehat P)h_\lambda\bigr\|_2^2
 =
 |a_\lambda|^2(1-\chi_0^2)
 +
 |b_\lambda|^2(1-\chi_4^2).
\]
Moreover, since \(b_\lambda\neq0\) and
\[
    1-\chi_4^2=(1-\chi_4)(1+\chi_4)
               \sim 2(1-\chi_4),
\]
one obtains
\[
    \bigl\|(1-\widehat P)h_\lambda\bigr\|_2^2
    \asymp 1-\chi_4.
\]
\end{proposition}
\paragraph{Proof sketch.}
The prolate modes diagonalize the concentration operator.  Orthogonality
therefore kills all cross terms in the leakage norm, giving the exact
quadratic sum above.  The \(S^3\) geometry identifies which modes occur and,
crucially, forces a nonzero \(h_4\)-component.  Since the leakage associated
to lower modes is smaller, the \(h_4\) term determines the asymptotic scale.
\paragraph{Relation to the Connes gap.}
CCM numerically observe a distinguished near-null direction in the quadratic
form.  The proposition gives a structural explanation: the small direction
is not an accidental numerical cancellation but the leakage of a specific
geometrically forced mode.  The near-kernel therefore has size
\[
    \Theta(1-\chi_4).
\]
This replaces an unexplained numerical phenomenon by a precise spectral
mechanism.
\subsection{Identification of the relevant obstruction scale}
\begin{corollary}[Single-scale reduction]
For the distinguished direction \(h_\lambda\), all three quantities
\[
    \widehat h_\lambda(0),\qquad
    h_\lambda(0),\qquad
    \|(1-\widehat P)h_\lambda\|_2^2
\]
are either zero or of order at most the common scale
\[
    \varepsilon_\lambda=1-\chi_4,
\]
with the leakage actually satisfying
\[
    \|(1-\widehat P)h_\lambda\|_2^2\asymp\varepsilon_\lambda.
\]
\end{corollary}
\paragraph{Proof sketch.}
Combine the preceding propositions.
\paragraph{Relation to the Connes gap.}
This is conceptually stronger than three unrelated estimates.  It shows that
the problematic direction in the compressed Weil form has a single small
parameter.  Thus the remaining question is no longer whether uncontrolled
errors from several sources can overwhelm positivity; one must compare the
Weil-form contribution against one explicitly identified tunnelling scale.
\subsection{Reduction of the remaining positivity problem to a sharp
trace-formula estimate}
Let \(Q\) denote the Weil quadratic form, \(W_\lambda\) the relevant
finite/prolate operator, and \(E h_\lambda\) the lifted geometric test vector.
\begin{proposition}[Reduction to the critical quadratic estimate]
After the exact cancellation and leakage estimates above, the unresolved
estimate in the distinguished near-null direction can be isolated as
\[
    QW_\lambda(Eh_\lambda)
    \asymp 1-\chi_4.
\]
Consequently the principal remaining analytic problem is to extract the
leading \(1-\chi_4\) contribution from the CCM trace formula.
\end{proposition}
\paragraph{Proof sketch.}
The geometric analysis removes the zero Fourier scalar exactly and bounds
the surviving endpoint scalar by \(O(1-\chi_4)\).  The leakage calculation
shows that the norm of the escaping component itself is
\(\Theta(1-\chi_4)\).  Therefore, once the already-controlled bulk terms are
separated, no larger unidentified error remains in this direction.
The question becomes whether the arithmetic/trace-formula contribution has a
nonvanishing leading coefficient on precisely this scale.
\paragraph{Status.}
This is a \emph{reduction}, not yet the final estimate:
\[
    \boxed{
    \widehat h_\lambda(0)=0,\qquad
    h_\lambda(0)=O(1-\chi_4)
    }
\]
are established geometrically, whereas
\[
    \boxed{
    QW_\lambda(Eh_\lambda)\asymp1-\chi_4
    }
\]
is the remaining analytic lemma to be extracted from the trace formula.
\paragraph{Relation to the Connes gap.}
This appears to isolate a concrete subproblem inside the general positivity
obstruction.  Rather than proving global positivity of the Weil form at once,
one needs to show that the arithmetic quadratic contribution does not vanish
to higher order than the geometrically determined tunnelling defect.
\subsection{Geometric interpretation of the CCM near-null vector}
\begin{proposition}[Boundary/tunnelling interpretation]
The CCM near-null direction is naturally interpreted as a boundary or
tunnelling mode of the \(S^3\)-induced prolate decomposition.  In particular,
its smallness is governed by
\[
    1-\chi_4,
\]
rather than by a generic approximation parameter.
\end{proposition}
\paragraph{Proof sketch.}
The preceding leakage identity identifies the dominant escaping component
with the fourth geometrically allowed mode.  The small eigenvalue defect
\(1-\chi_4\) measures precisely the failure of that mode to remain inside the
concentrated spectral region.  Hence the small quadratic direction is tied
to spectral tunnelling across the finite cutoff.
\paragraph{Relation to the Connes gap.}
One of the difficulties in the CCM program is that the numerically observed
small eigenvalue could in principle signal a genuine failure of positivity.
The \(S^3\) picture instead supplies a mechanism predicting both its existence
and its scale.  Thus the small eigenvalue itself ceases to be mysterious;
what remains is to determine its sign and leading arithmetic coefficient.
\subsection{Net effect on the Connes program}
The concrete contribution of the \(S^3\) construction can therefore be
summarized as the chain
\[
\boxed{
\begin{array}{c}
\text{unstructured finite-dimensional positivity defect}
\\[2mm]
\Downarrow\quad S^3\text{ mode decomposition}
\\[2mm]
\widehat h_\lambda(0)=0,
\qquad
h_\lambda(0)=O(1-\chi_4)
\\[2mm]
\Downarrow
\\[2mm]
\|(1-\widehat P)h_\lambda\|^2
   \asymp1-\chi_4
\\[2mm]
\Downarrow
\\[2mm]
\text{one critical tunnelling direction}
\\[2mm]
\Downarrow
\\[2mm]
QW_\lambda(Eh_\lambda)
   \stackrel{?}{\asymp}1-\chi_4 .
\end{array}}
\]
Thus the strongest current claim is not that the \(S^3\) geometry proves the
Connes positivity criterion.  It is that it has converted part of the
previously diffuse obstruction into a one-dimensional, single-scale
asymptotic problem whose missing input is now located on the arithmetic side
of the trace formula.
In particular, the genuinely new reductions are:
\[
\begin{array}{lll}
\text{(i)}&
\widehat h_\lambda(0)=0,
&
\text{exact removal of an exceptional scalar mode},
\\[1mm]
\text{(ii)}&
h_\lambda(0)=O(1-\chi_4),
&
\text{endpoint obstruction reduced to tunnelling scale},
\\[1mm]
\text{(iii)}&
\|(1-\widehat P)h_\lambda\|^2\asymp1-\chi_4,
&
\text{rigorous scale of the CCM near-null direction},
\\[1mm]
\text{(iv)}&
\text{positivity defect }\rightsquigarrow
QW_\lambda(Eh_\lambda)\asymp1-\chi_4,
&
\text{isolation of the remaining trace-formula lemma}.
\end{array}
\]
Items (i)--(iii) are geometric/spectral improvements already obtained.
Item (iv) is the resulting reduction and is currently the point at which the
\(S^3\) argument hands the problem back to the arithmetic trace formula.
